\documentclass[10pt,a4paper]{amsart}
\usepackage[margin=19mm]{geometry}
\usepackage{amsmath,amssymb,mathtools}
\usepackage[scr=rsfs,cal=euler]{mathalfa}
\usepackage{xfrac}
\usepackage[unicode,hidelinks,bookmarks=false]{hyperref}
\newcommand{\ie}{i.e.\ }
\newcommand{\cf}{cf.\ }
\newcommand{\resp}{respectively\ }
\newcommand{\Subsection}{\subsection}
\providecommand{\nobreakdash}{\nobreak\hbox{-}}
\setlength{\emergencystretch}{2em}
\multlinegap0pt
\usepackage{amssymb}
\usepackage{enumerate}
\usepackage{tikz}
\usepackage{mdframed}
\newtheorem{lemma}{Lemma}
\newcommand{\Ga}{\Gamma}
\newcommand{\La}{\Lambda}
\newcommand{\R}{\mathbb{R}}
\newcommand{\Z}{\mathbb{Z}}
\newcommand{\bB}{\mathbf{B}}
\newcommand{\bE}{\mathbf{E}}
\newcommand{\cB}{\mathcal{B}}
\newcommand{\cF}{\mathcal{F}}
\newcommand{\cH}{\mathcal{H}}
\newcommand{\cL}{\mathcal{L}}
\newcommand{\cM}{\mathcal{M}}
\newcommand{\cN}{\mathcal{N}}
\newcommand{\cP}{\mathcal{P}}
\newcommand{\cQ}{\mathcal{Q}}
\newcommand{\cR}{\mathcal{R}}
\newcommand{\cT}{\mathcal{T}}
\newcommand{\cV}{\mathcal{V}}
\newcommand{\e}{\varepsilon}
\renewcommand{\geq}{\geqslant}
\renewcommand{\leq}{\leqslant}
\renewcommand{\vec}[1]{\mathbf{#1}}
\title{On the spectrum of Diophantine exponents of lattices}
\author{Oleg\, N.\, German}
\date{Source: arXiv:2603.06016v1 (CC BY 4.0)}
\begin{document}
\maketitle

\noindent\emph{Source and licence.} On the spectrum of Diophantine exponents of lattices. Version arXiv:2603.06016v1 (CC BY 4.0), distributed by arXiv under the Creative Commons Attribution 4.0 International licence, http://creativecommons.org/licenses/by/4.0/, as recorded in the arXiv licence metadata for this exact version and shown on the abstract page https://arxiv.org/abs/2603.06016v1. Full licence notice: \texttt{licenses/arxiv-2603.06016-CC-BY-4.0.txt}.\par\smallskip

\noindent\textit{Editorial scope.} Counted unit: the article's lemma l:main\_metric\_argument (source0.tex lines 724--734). The article's complete original proof of it is delivered with it (source0.tex lines 736--1063, one \texttt{proof} environment of 328 lines). No mathematics is written, repaired or re-derived: the statement and the proof are the article's own lines, in the article's own notation, and no retained line is substituted or re-flowed.  No further source-local context is needed: every cross-reference of every retained line resolves inside the two delivered spans.  Nothing was omitted from any retained span, so the package carries no omission record.  No retained line contains a citation, so the wrapper carries no bibliography and no bibliography entry.  No retained line contains a figure environment, an image-inclusion command or drawing code, so no image is delivered and the package declares no asset dependency.  Editorial formatting only, in addition: an AMS \texttt{amsart} preamble that restates the article's own declarations and macros used by the retained lines, namely the environments \texttt{lemma} on separate counters, together with the standard packages those lines need.  The retained environments print in this document as Lemma 1 in order of appearance, whereas the article prints the same items with the numbering of its own full document.  The complete native source of the article as distributed in the arXiv e-print for this exact version is bundled unchanged as \texttt{sources/195834/source0.tex}.
\begin{lemma}\label{l:main_metric_argument}
  Let the cube\,\ $\bB$ be chosen so that for each $\bE=(\vec e_1,\ldots,\vec e_n)\in\bB$ the vectors $\vec v_1,\vec v_2,\vec e_1,\ldots,\vec e_n$ are linearly independent. Let us set for each $\bE\in\bB$
  \[
    \La_\bE=\Z\vec e_1+\ldots+\Z\vec e_n+\Z\vec v_1+\Z\vec v_2.
  \]
  Then for every $\e>0$ there is a constant $c=c(\e)>0$ such that for almost every $\bE\in\bB$ the inequality
  \[
    \Pi(\vec x)>\frac c{\log^{1+\e}(1+|\vec x|)}
  \]
  holds for every $\vec x\in\La_\bE\backslash\Ga_\cL$.
\end{lemma}

\begin{proof}
  Let us fix $\e>0$. For all positive $t$, set
  \[
    \varphi(t)=\frac1{\log^{d+\e}(1+t)}\,.
  \]
  Let us define the ``forbidden'' set
  \[
    \cF=\cF(\varphi)=\Big\{\vec x\in\R^d \,\Big|\, \Pi(\vec x)^d\leq\varphi(|\vec x|),\ |\vec x|\geq1 \Big\}
  \]
  and consider the ``exceptional'' set
  \[
    \cM=\Big\{\bE\in\bB \,\Big|\, \text{the set }\big(\La_\bE\backslash\Ga_\cL\big)\cap\cF\text{ is infinite} \Big\}.
  \]
  Let us show that $\mu(\cM)=0$, where $\mu$ is the Lebesgue measure on the space $\R^{nd}$. We split our argument into several steps.
  
  \paragraph{Reduction to Borel--Cantelli lemma.}
  
  Each vector $\vec x\in\La_\bE$ enjoys a representation
  \begin{equation}\label{eq:a_point_of_La_E}
    \vec x=z_1\vec e_1+\ldots+z_n\vec e_n+z_{n+1}\vec v_1+z_{n+2}\vec v_2
  \end{equation}
  with integer $z_1,\ldots,z_{n+2}$ (recall that $n=d-2$). And conversely, for a fixed $\bE$, each $\vec z=(z_1,\ldots,z_d)\in\Z^d$ determines a vector $\vec x=\vec x(\vec z,\bE)$ of the form \eqref{eq:a_point_of_La_E}. Moreover,
  \[
    \vec x(\vec z,\bE)\in\Ga_\cL
    \iff
    \vec z\in\cV,
  \]
  where $\cV$ is the coordinate subspace of the last two coordinates.

  For each $\nu\in\Z$, $\nu\geq0$, let us consider the set
  \[
    \cM_\nu=\Big\{\bE\in\bB \,\Big|\, \exists\,\vec z\in\Z^d\backslash\cV:\,2^\nu\leq|\vec z|<2^{\nu+1},\ \vec x(\vec z,\bE)\in\cF \Big\}.
  \]
  Then
  \[
    \cM=\bigcap_{\nu_0=0}^\infty\bigcup_{\nu=\nu_0}^\infty\cM_\nu.
  \]
  By Borel--Cantelli lemma it suffices to prove the convergence of the series
  \begin{equation}\label{eq:borel_cantelli_series}
    \sum_{\nu=0}^\infty\mu(\cM_\nu).
  \end{equation}
  
  \paragraph{Sets $\cF^i_\nu$ and $\cM^i_\nu$.}  
  
  Note that if
  \[
    2^\nu\leq|\vec z|<2^{\nu+1},
  \]
  then by compactness of $\bB$ there is a positive integer $\eta$ depending only on $\bB$ such that
  \[
    2^{\nu-\eta}\leq|\vec x(\vec z,\bE)|\leq2^{\nu+\eta}.
  \]
  Set
  \[
    \cF_\nu=\Big\{\vec x\in\R^d \,\Big|\, \Pi(\vec x)^d\leq\varphi(|\vec x|),\ 2^{\nu-\eta}\leq|\vec x|\leq2^{\nu+\eta} \Big\}.
  \]
  Set also, for each $i=1,\ldots,d$,
  \[
    \cF^i_\nu=\Big\{\vec x\in\cF_\nu \,\Big|\, |\vec x|=x_i \Big\}.
  \]
  Then
  \[
    \cM_\nu\subseteq\bigcup_{i=1}^d\cM^i_\nu,
  \]
  where
  \[
    \cM^i_\nu=\Big\{\bE\in\bB \,\Big|\, \exists\,\vec z\in\Z^d\backslash\cV:\,2^\nu\leq|\vec z|<2^{\nu+1},\ \vec x(\vec z,\bE)\in\cF^i_\nu \Big\}.
  \]
  
  \paragraph{Covering $\cF^1_\nu$ by parallelepipeds.}
  
  Let us cover the set $\cF^1_\nu$ by a family of paralle\-le\-pi\-peds. This will allow us to estimate the measure of $\cM^1_\nu$\,. The measure of all the other $\cM^i_\nu$ is estimated in the same way.

  For each $\vec x=(x_1,x_2,\ldots,x_d)\in\R^d$, set $\underline{\vec x}=(x_2,\ldots,x_d)\in\R^{d-1}$.
  For each $\tau\in\Z$, $\tau\geq0$, define the $(d-1)$-dimensional set
  \[
    \cH_\tau=\Big\{\underline{\vec x}=(x_2,\ldots,x_d)\in\R^{d-1} \,\Big|\, |\underline{\vec x}|\leq2^{\tau+1},\ 
                                                                            %\Pi(\underline{\vec x})^{d-1}
                                                                            |x_2\cdot\ldots\cdot x_d|\leq\frac{\varphi(2^\tau)}{2^\tau} \Big\}.
  \]
  Then 
  \begin{equation}\label{eq:F_1_nu_hyperbolic_covering}
    \cF^1_\nu\subset\bigcup_{\tau=\nu-\eta}^{\nu+\eta}\Big([2^\tau,2^{\tau+1}]\times\cH_\tau\Big).
  \end{equation}

  Let us cover the set $\cH_\tau$ by parallelepipeds as follows. Define an integer $\sigma=\sigma(\tau)$ by the inequalities
  \[
    2^{-\sigma-1}<
    \bigg(\dfrac{\varphi(2^\tau)}{2^\tau}\bigg)^{\frac1{d-1}}\leq
    2^{-\sigma}.
  \]
  Then
  \begin{equation}\label{eq:sigma_of_tau}
    \sigma=\sigma(\tau)=\bigg[\frac{\tau-\log_2(\varphi(2^\tau))}{d-1}\bigg].
  \end{equation}
  Consider the family of $(d-1)$-tuples of integers
  \[
    \cT=\cT(\tau)=\Big\{\underline{\pmb\tau}=(\tau_2,\ldots,\tau_d)\in\Z^{d-1} \,\Big|\, \tau_2+\ldots+\tau_d=0,\ |\underline{\pmb\tau}|\leq(d-1)(\sigma+\tau+2) \Big\}
  \]  
  and determine for each $\underline{\pmb\tau}\in\cT$ the parallelepiped
  \[
    \cP(\underline{\pmb\tau})=
    \Big\{\underline{\vec x}=(x_2,\ldots,x_d)\in\R^{d-1} \,\Big|\, |x_i|\leq2^{\tau_i-\sigma+(d-2)},\ i=2,\ldots,d \Big\}.
  \]
  Let us show that
  \begin{equation}\label{eq:H_tau_covering}
    \cH_\tau\subset\bigcup_{\underline{\pmb\tau}\in\cT}\cP(\underline{\pmb\tau}).
  \end{equation}
  Indeed, let $\underline{\vec y}\in\cH_\tau$\,, $|y_2\cdot\ldots\cdot y_d|=\dfrac{\varphi(2^\tau)}{2^\tau}$\,. Set
  \[
    \underline{\vec w}=(w_2,\ldots,w_d)=
    \underline{\vec y}\cdot2^{-\sigma}\cdot\bigg(\dfrac{\varphi(2^\tau)}{2^\tau}\bigg)^{-\frac1{d-1}}.
  \]
  The points $\underline{\vec y}$ and $\underline{\vec w}$ are proportional,
  \begin{equation}\label{eq:bounds_for_w}
    |\underline{\vec y}|\leq|\underline{\vec w}|<2|\underline{\vec y}|\leq2^{\tau+2},
    \qquad\text{ and }\qquad
    |w_2\cdot\ldots\cdot w_d|=2^{-(d-1)\sigma}.
  \end{equation}
  Set $\tau_i=\sigma+\big\lceil\log_2|w_i|\big\rceil$, $i=2,\ldots,d-1$ and $\tau_d=-\tau_2-\ldots-\tau_{d-1}$. Here $\lceil\,\cdot\,\rceil$ denotes the ceiling function. Then by \eqref{eq:bounds_for_w}
  \[
    \quad
    2^{\tau_i-\sigma-1}<|w_i|\leq2^{\tau_i-\sigma},
    \qquad\qquad\ \
    |\tau_i|\leq\sigma+\tau+2,
    \qquad
    i=1,\ldots,d-1,
  \]
  \[
    |w_d|<
    2^{\tau_d-\sigma+(d-2)},
    \qquad
    |\tau_d|\leq(d-1)(\sigma+\tau+2).
  \]
  Thus, $\underline{\vec w}$ and, therefore, $\underline{\vec y}$ belong to the parallelepiped $\cP(\underline{\pmb\tau})$. I.e., indeed, inclusion \eqref{eq:H_tau_covering} takes place.
  
  Taking into account \eqref{eq:F_1_nu_hyperbolic_covering}, we get the following covering of $\cF^1_\nu$ by parallelepipeds:
  \[
    \cF^1_\nu=
    \bigcup_{\substack{\pmb\tau=(\tau_1,\ldots,\tau_d)\,\in\,\Z^d \\ 
             \nu-\eta\,\leq\,\tau_1\,\leq\,\nu+\eta \vphantom{\frac11} \\ 
             \underline{\pmb\tau}\,\in\,\cT(\tau_1)}}
    \cQ(\pmb\tau),
    \qquad\text{ where }\qquad
    \cQ(\pmb\tau)=
    [2^{\tau_1},2^{\tau_1+1}]\times\cP(\underline{\pmb\tau}).
  \]
  
  This covering in turn gives the following covering of the set $\cM^1_\nu$:
  \[
    \cM^1_\nu=
    \bigcup_{\substack{\pmb\tau=(\tau_1,\ldots,\tau_d)\,\in\,\Z^d \\ 
             \nu-\eta\,\leq\,\tau_1\,\leq\,\nu+\eta \vphantom{\frac11} \\ 
             \underline{\pmb\tau}\,\in\,\cT(\tau_1)}}
    \bigcup_{\substack{\vec z\in\Z^d\backslash\cV \\ 
             \,\ 2^\nu\leq|\vec z|<2^{\nu+1}}}
    \cN(\vec z,\pmb\tau),
  \]
  where
  \[
    \cN(\vec z,\pmb\tau)=
    \Big\{\bE\in\bB \,\Big|\, \vec x(\vec z,\bE)\in\cQ(\pmb\tau) \Big\}.
  \]
  Thus, the measure of $\cM^1_\nu$ can be estimated as
  \begin{equation}\label{eq:estimate_by_measure_of_N}
    \mu(\cM^1_\nu)\leq
    \sum_{\tau_1=\nu-\eta}^{\nu+\eta}\phantom{1}
    \sum_{\underline{\pmb\tau}\,\in\,\cT(\tau_1)}
    \sum_{\substack{\vec z\in\Z^d\backslash\cV \\ 
                    \,\ 2^\nu\leq|\vec z|<2^{\nu+1}}}
    \mu(\cN(\vec z,\pmb\tau)).
  \end{equation}
  
  \paragraph{\bf Analysis of $\cN(\vec z,\pmb\tau)$}
  
  Let us write the coordinates of $\vec e_1,\ldots,\vec e_n$ and $\vec v_1,\vec v_2$ in the columns of matrices $E$ and $V$, respectively:
  \[
    E=
    \begin{pmatrix}
      e_{1,1} & \cdots & e_{n,1} \\
      \vdots  & \ddots & \vdots \\
      e_{1,d} & \cdots & e_{n,d}
    \end{pmatrix},
    \qquad
    V=
    \begin{pmatrix}
      v_{1,1} & v_{2,1} \\
      \vdots  & \vdots \\
      v_{1,d} & v_{2,d}
    \end{pmatrix}.
  \]
  Denote by $\vec e^1,\ldots,\vec e^d$ the rows of $E$ and by $\vec v^1,\ldots,\vec v^d$ -- the rows of $V$. Let us set for $\vec z=(z_1,\ldots,z_d)$
  \[
    \vec z_E=(z_1,\ldots,z_n)
    \qquad\text{ and }\qquad
    \vec z_V=(z_{n+1},z_{n+2}).
  \]
  Let us denote by $\langle\,\cdot\,,\,\cdot\,\rangle$ the inner product. Then the condition $\vec x(\vec z,\bE)\in\cQ(\pmb\tau)$ is equivalent to the system of inequalities
  \begin{equation}\label{eq:forbidden_set_system}
    \begin{cases}
      2^{\tau_1}\leq\langle\vec e^1,\vec z_E\rangle+\langle\vec v^1,\vec z_V\rangle\leq2^{\tau_1+1} \\
      |\langle\vec e^2,\vec z_E\rangle+\langle\vec v^2,\vec z_V\rangle|\leq2^{\tau_2-\sigma+(d-2)} \\
      \ \cdot\cdot\cdot\cdot\cdot\cdot\cdot\cdot\cdot\cdot\cdot\cdot\cdot\cdot\cdot\cdot\cdot\cdot\cdot\cdot\cdot\cdot\cdot\cdot\cdot\cdot\cdot \\
      |\langle\vec e^d,\vec z_E\rangle+\langle\vec v^d,\vec z_V\rangle|\leq2^{\tau_d-\sigma+(d-2)}
    \end{cases},
  \end{equation}
  where $\sigma=\sigma(\tau_1)$ is defined by \eqref{eq:sigma_of_tau}.
  
  For every $j=1,\ldots,d$, the set of all possible values of $\vec e^j$, i.e. the set
  \[
    \Big\{\vec e^j \,\Big|\, \bE\in\bB \Big\},
  \]
  is a unit cube in $\R^n$. Denote this cube by $\cB^j$. Then
  \[
    \cN(\vec z,\pmb\tau)=
    \cN^1(\vec z,\pmb\tau)\times\ldots\times\cN^d(\vec z,\pmb\tau),
  \]
  where
  \[
    \cN^1(\vec z,\pmb\tau)=
    \Big\{\vec e^1\in\cB^1 \,\Big|\, 2^{\tau_1}\leq\langle\vec e^1,\vec z_E\rangle+\langle\vec v^1,\vec z_V\rangle\leq2^{\tau_1+1} \Big\}
  \]
  and
  \[
    \cN^j(\vec z,\pmb\tau)=
    \Big\{\vec e^j\in\cB^j \,\Big|\, |\langle\vec e^j,\vec z_E\rangle+\langle\vec v^j,\vec z_V\rangle|\leq2^{\tau_j-\sigma+(d-2)} \Big\}
  \]
  for $j=2,\ldots,d$.
  
  We estimate roughly the measure of $\cN^1(\vec z,\pmb\tau)$:
  \[
    \mu(\cN^1(\vec z,\pmb\tau))\leq1.
  \]
  
  The set $\cN^j(\vec z,\pmb\tau)$ for a given $j=2,\ldots,d$, in general, can be empty. If it is not empty, its measure can be obviously estimated as
  \[
    \mu(\cN^j(\vec z,\pmb\tau))\ll
    \frac{2^{\tau_j-\sigma+(d-2)}}{|\vec z_E|}\ll_d
    \frac{2^{\tau_j-\sigma}}{|\vec z_E|}\,.
  \]
  Since $\tau_2+\ldots+\tau_d=0$ and $\sigma=\sigma(\tau_1)$ is defined by \eqref{eq:sigma_of_tau}, we get
  \begin{equation}\label{eq:estimating_measure_of_N}
    \mu(\cN(\vec z,\pmb\tau))\ll_d
    \frac{2^{-(d-1)\sigma}}{|\vec z_E|^{d-1}}\ll_d
    \frac{2^{-\tau_1+\log_2(\varphi(2^{\tau_1}))}}{|\vec z_E|^{d-1}}=
    \frac{2^{-\tau_1}\cdot\varphi(2^{\tau_1})}{|\vec z_E|^{d-1}}\,.
  \end{equation}
  Note that this estimate is uniform in $\underline{\pmb\tau}$\,. And it is nontrivial only in the case when each of the sets $\cN^j(\vec z,\pmb\tau)$, $j=2,\ldots,d$, is not empty.
  
  \paragraph{Estimating the number of nonempty $\cN^j(\vec z,\pmb\tau)$ for a fixed $\vec z_E$.}
  
  Recall that the vectors $\vec v^j$, $j=2,\ldots,d$, are fixed and distinct from the zero vector (as the coordinates of $\vec v_1$, $\vec v_2$ are positive). Let us fix a vector $\vec z_E$, an index $j$, $2\leq j\leq d$, and a positive constant $C$. Let us ask ourselves what is the cardinality of the set
  \[
    \cR^j(\vec z_E,C)=
    \Big\{\vec z_V\in\Z^2 \,\Big|\, |\vec z_V|\leq C,\ \exists\,\vec e\in\cB^j:\ |\langle\vec e,\vec z_E\rangle+\langle\vec v^j,\vec z_V\rangle|\leq1 \Big\}.
  \]
  The set of values of $\langle\vec e,\vec z_E\rangle$ as $\vec e$ ranges through the cube $\cB^j$ is a segment of length of order $|\vec z_E|$. Hence the set $\cR^j(\vec z_E,C)$ is contained in a rectangle with sides of order $C$ and $|\vec z_E|$. The number of integer points in such a rectangle is bounded by a quantity of the order of its area, i.e.
  \[
    |\cR^j(\vec z_E,C)|\ll C|\vec z_E|.
  \]
  Finally, we note that since
  \[
    \sum_{j=2}^{d}\big(\tau_j-\sigma+(d-2)\big)=-(d-1)\sigma+O(d)<-\tau_1+\log_2(\varphi(2^{\tau_1}))+O(d),
  \]
  and $\tau_1$ is large, there is $j$ such that $2^{\tau_j-\sigma+(d-2)}<1$. Hence, for the set
  \[
    \cR(\vec z_E,C)=
    \bigcap_{j=2}^d\cR^j(\vec z_E,C),
  \]
  we have
  \begin{equation}\label{eq:estimate_on_R}
    |\cR(\vec z_E,C)|\ll_\bB C|\vec z_E|.
  \end{equation}
  
  \paragraph{Estimating the measure of $\cM_\nu$.}
  
  By \eqref{eq:estimate_by_measure_of_N}, \eqref{eq:estimating_measure_of_N}, \eqref{eq:estimate_on_R} we get the following estimate for the measure of $\cM^1_\nu$\,:
  \begin{align*}
    \mu(\cM^1_\nu)\, & \leq
    \sum_{\tau_1=\nu-\eta}^{\nu+\eta}\phantom{1}
    \sum_{\underline{\pmb\tau}\,\in\,\cT(\tau_1)}
    \sum_{\substack{\vec z\in\Z^d\backslash\cV \\ 
                    \,\ 2^\nu\leq|\vec z|<2^{\nu+1}}}
    \mu(\cN(\vec z,\pmb\tau))
    \ll_\bB \\ & \ll_\bB
    \sum_{\tau_1=\nu-\eta}^{\nu+\eta}\phantom{1}
    \sum_{\underline{\pmb\tau}\,\in\,\cT(\tau_1)}
    \sum_{\substack{\vec z_E\in\Z^n \\ 
                    \,\ 0<|\vec z_E|<2^{\nu+1}}}
    \sum_{\vec z_V\in\cR(\vec z_E,2^{\nu+1})}
    \frac{2^{-\tau_1}\cdot\varphi(2^{\tau_1})}{|\vec z_E|^{d-1}}
    \ll_\bB \\ & \ll_\bB
    \sum_{\tau_1=\nu-\eta}^{\nu+\eta}
    |\cT(\tau_1)|\cdot2^{-\tau_1}\cdot\varphi(2^{\tau_1})
    \hskip-5mm
    \sum_{\substack{\vec z_E\in\Z^n \\ 
                    \,\ 0<|\vec z_E|<2^{\nu+1}}}
    \frac{2^{\nu+1}}{|\vec z_E|^{d-2}}
    \ll_\bB \\ & \ll_\bB
    \sum_{\tau_1=\nu-\eta}^{\nu+\eta}
    \big(d\tau_1-\log_2(\varphi(2^{\tau_1})\big)^{d-2}\cdot2^{-\tau_1+\nu+1}\cdot\varphi(2^{\tau_1})
    \sum_{T=1}^{2^{\nu+1}}
    \sum_{\substack{\vec z_E\in\Z^n \\ 
                    |\vec z_E|=T}}
    \frac1{|\vec z_E|^{d-2}}
    \ll_\bB \\ & \ll_\bB
    \Big(\nu+\big|\log_2\big(\varphi(2^{\nu+\eta})\big)\big|\Big)^{d-2}\hskip-1mm\cdot\varphi(2^{\nu-\eta})
    \sum_{T=1}^{2^{\nu+1}}\frac1T
    \ll_\bB \\ & \ll_\bB
    \ \nu\cdot\varphi(2^{\nu-\eta})\cdot\Big(\nu^{d-2}+\big|\log_2(\varphi\big(2^{\nu+\eta})\big)\big|^{d-2}\Big).
    \vphantom{\frac{\Big|}{}}
  \end{align*}
  
  The measure of the sets $\cM^2_\nu,\ldots,\cM^d_\nu$ is estimated in the same way. Hence
  \[
    \mu(\cM_\nu)\,\ll_\bB\,
    \varphi(2^{\nu-\eta})\cdot\Big(\nu^{d-1}+\nu\big|\log_2(\varphi\big(2^{\nu+\eta})\big)\big|^{d-2}\Big),
  \]
  where $\eta$ is a positive constant depending on $\bB$. Taking into account that
  \[
    \varphi(t)=\frac1{\log^{d+\e}(1+t)}\,,
  \]
  we get
  \[
    \mu(\cM_\nu)\ll_\bB\nu^{-1-\e},
  \]
  i.e. the series \eqref{eq:borel_cantelli_series} converges. It follows by Borel--Cantelli lemma that the measure of the ``exceptional'' set $\cM$ is equal to zero. Thus, for almost every $\bE\in\bB$ the inequality $\Pi(\vec x)^d\leq\varphi(|\vec x|)$ holds only for a finite number of points of the set $\La_\bE\backslash\Ga_\cL$. Which immediately implies the desired statement.
\end{proof}

\end{document}
